\section*{Hoofdstuk \ref{ch:b1:stress-strain-rheology} --- Spanning, rek en reologie}

\begin{solution}{exo:b1:stress-strain-rheology:1}
(a) $\sigma_v = 2700\times9.81\times10\,000 = \qty{2.65e8}{Pa} = \qty{264.9}{MPa}$.
(b) $P = 1000\times9.81\times10\,000 = \qty{98.1}{MPa}$. (c) $\sigma_v' =
264.9 - 98.1 = \qty{166.8}{MPa}$, en $\lambda = 98.1/264.9 = 0.370$, wat
gelijk is aan $1000/2700$.
\end{solution}

\begin{solution}{exo:b1:stress-strain-rheology:2}
(a) $e = (13-10)/10 = 0.30$ en $S = 1.30$. (b) $e = (7-10)/10 = -0.30$, een
verkorting van \qty{30}{\percent}, en $S = 0.70$. (c) De hoek is met
$\psi = \ang{30}$ gesloten, en $\gamma = \tan\ang{30} = 0.577$.
\end{solution}

\begin{solution}{exo:b1:stress-strain-rheology:3}
(a) Middelpunt \qty{70}{MPa}, straal \qty{30}{MPa}. (b) De grootste schuifspanning
is de straal, \qty{30}{MPa}, op de vlakken onder \ang{45} met $\sigma_1$ en
$\sigma_3$, die de \omterm{def:b1:stress-strain-rheology:stress}{normaalspanning} van het middelpunt dragen, \qty{70}{MPa}. (c)
$2\theta = \ang{60}$: $\sigma_n = 70 + 30\cos\ang{60} = \qty{85}{MPa}$ en
$\tau = 30\sin\ang{60} = \qty{26.0}{MPa}$.
\end{solution}

\begin{solution}{exo:b1:stress-strain-rheology:4}
(a) $\ang{45} - \ang{30}/2 = \ang{30}$ tot $\sigma_1$, aan weerszijden ervan,
zodat de twee vlakken \ang{60} met elkaar maken. (b) Met $c' = 20 + r$ leest de
voorwaarde van \cref{thm:b1:stress-strain-rheology:failure-angle} als
$r = C\cos\phi + (20 + r)\sin\phi$, dus
\[
  r = \frac{C\cos\phi + 20\sin\phi}{1 - \sin\phi} = \frac{8.66 + 10}{0.5}
    = \qty{37.3}{MPa},
\]
en $\sigma_1' = 20 + 2r = \qty{94.6}{MPa}$.
\end{solution}

\begin{solution}{exo:b1:stress-strain-rheology:5}
(a) $c = 40$ en $R = \sqrt{20^2 + 15^2} = \qty{25}{MPa}$: $\sigma_1 =
\qty{65}{MPa}$ en $\sigma_3 = \qty{15}{MPa}$; $\tan2\theta_p = 30/40 = 0.75$,
$2\theta_p = \ang{36.87}$, $\theta_p = \ang{18.43}$. (b) Bij deze hoek zijn
$\cos^2\theta = 0.9$, $\sin^2\theta = 0.1$ en $\sin\theta\cos\theta = 0.3$:
$\sigma_n = 60\times0.9 + 2\times15\times0.3 + 20\times0.1 = 54 + 9 + 2 =
\qty{65}{MPa}$, en $\tau = (20-60)\times0.3 + 15\times(0.9-0.1) = -12 + 12 =
0$. (c) $R = \qty{25}{MPa}$, op de vlakken onder \ang{45} met de hoofdrichtingen.
\end{solution}

\begin{solution}{exo:b1:stress-strain-rheology:6}
(a) $S_{1,2} = \sqrt{1 + 1} \pm 1$: $S_1 = 2.414$ en $S_2 = 0.414$;
$\tan2\theta' = 2/2 = 1$, $\theta' = \ang{22.5}$. (b) $S_1S_2 =
(\sqrt2 + 1)(\sqrt2 - 1) = 1$, wat gelijk is aan $\det F$: geen verandering van
oppervlak. (c) De lijn langs $(1, 1)$ gaat naar $F(1, 1) = (1 + 2, 1) = (3, 1)$,
onder $\arctan(1/3) = \ang{18.4}$ met de schuifrichting, met de \omterm{def:b1:stress-strain-rheology:strain}{strekking}
$\sqrt{10}/\sqrt2 = 2.24$. Zij is voorbij de lange as van de ellips gedraaid,
bij \ang{22.5}, richting het schuifvlak; haar \omterm{def:b1:stress-strain-rheology:strain}{strekking} is iets kleiner dan
$S_1$, zoals het moet, want geen lijn strekt meer dan $S_1$.
\end{solution}

\begin{solution}{exo:b1:stress-strain-rheology:7}
(a) $S_1/S_2 = 4$ en $S_1S_2 = 1$ geven $S_1 = 2$ en $S_2 = 0.5$; volgens
\cref{prop:b1:stress-strain-rheology:simple-shear} is $\gamma = S_1 - S_2 =
1.5$. (b) $\tan2\theta' = 2/1.5 = 1.333$, $2\theta' = \ang{53.1}$: de lange assen
moeten onder \ang{26.6} met de wanden liggen. (c) $1.5\times200 = \qty{300}{m}$.
\end{solution}

\begin{solution}{exo:b1:stress-strain-rheology:8}
$\sigma_v = 2700\times9.81\times3000 = \qty{79.46}{MPa}$. (a) $P = \qty{29.43}{MPa}$,
$\sigma_n' = 70 - 29.43 = \qty{40.57}{MPa}$, en Byerlee staat
$0.85\times40.57 = \qty{34.5}{MPa}$ toe, meer dan de \qty{30}{MPa} die wordt aangelegd: de
breuk houdt stand. (b) $P = 0.6\times79.46 = \qty{47.68}{MPa}$, $\sigma_n' =
\qty{22.32}{MPa}$, en de wrijving geeft slechts \qty{19.0}{MPa}: de breuk
glijdt. (c) $30 = 0.85\,(70 - 79.46\,\lambda)$ geeft $\lambda = (70 -
35.29)/79.46 = 0.437$.
\end{solution}

\begin{solution}{exo:b1:stress-strain-rheology:9}
(a) De spanning gaat als $\dot\varepsilon^{1/4}$: $10^{1/4} = 1.78$. (b) De
verhouding van de sterkten is
\[
  \exp\Bigl[\frac{E}{nR}\Bigl(\frac{1}{673.15} - \frac{1}{773.15}\Bigr)\Bigr]
  = \exp\bigl(6705\times\num{1.921e-4}\bigr) = \exp(1.288) = 3.63 .
\]
(c) Een tienvoudige verandering van snelheid verandert de kruipsterkte met een
factor 1.78, wat de temperatuur in slechts $100\times\ln1.78/\ln3.63 \approx
\qty{45}{\degreeCelsius}$ doet; bij een \omterm{def:b1:heat-budget:geotherm}{geotherm} van ongeveer \qty{20}{\degreeCelsius}
per kilometer verplaatst dat de overgang met ongeveer \qty{2}{km}. Een verandering
van de \omterm{def:b1:heat-budget:geotherm}{geotherm} verplaatst de isothermen, en daarmee de overgang, veel meer.
\end{solution}

\begin{solution}{exo:b1:stress-strain-rheology:10}
(a) Met $C = 0$ glijdt een breuk wanneer $r = \mu c'/\sqrt{1+\mu^2}$
(\cref{thm:b1:stress-strain-rheology:failure-angle}). Hier is $c' = S =
\rho g z(1-\lambda)$, dus $\sigma_1 - \sigma_3 = 2r = 2\mu S/\sqrt{1+\mu^2}$.
(b) $S$ neemt toe met $2700\times9.81\times1000\times(1 - 1000/2700) =
\qty{16.68}{MPa}$ per kilometer, en $2\times0.85/1.3124 = 1.295$: de helling
is \qty{21.6}{MPa} per kilometer, tussen \num{13.1} (normaalbreuken) en
\qty{61.3}{MPa} per kilometer (overschuivingen). (c) $2\mu/(k+\mu) = 0.786$,
$2\mu/k = 1.295$ en $2\mu/(k-\mu) = 3.68$, met $k = 1.3124$: de
\omterm{def:b1:stress-strain-rheology:principal}{differentiële spanning} bij het glijden is 0.79, 1.30 en 3.68 keer de verticale
\omterm{def:b1:stress-strain-rheology:lithostatic}{effectieve spanning}.
\end{solution}

\begin{solution}{exo:b1:stress-strain-rheology:11}
(a) $t = 60\times\num{e6}\times\num{3.156e7} = \qty{1.89e15}{s}$ en
$2\sqrt{\kappa t} = 2\sqrt{\num{1.89e9}} = \qty{87.0}{km}$. (b)
$\operatorname{erf}x = 600/1350 = 0.444$ voor $x = 0.417$: $z = 0.417\times87.0
= \qty{36.3}{km}$; $\operatorname{erf}x = 800/1350 = 0.593$ voor $x = 0.586$:
$z = \qty{51.0}{km}$. (c) De aardbevingen zouden \qtyrange{36}{51}{km} moeten
bereiken; de omhullende legt de basis van de brosse laag op \qty{38.2}{km},
binnen dat bereik en dicht bij zijn top, bij een temperatuur, \qty{641}{\degreeCelsius},
binnen het waargenomen bereik.
\end{solution}

\begin{solution}{exo:b1:stress-strain-rheology:12}
(a) $r = 45$ en $c' = \qty{65}{MPa}$; de afstand van het middelpunt tot de lijn
$\tau = 0.85\sigma_n'$ is $0.85\times65/1.3124 = \qty{42.1}{MPa}$, kleiner dan
de straal: de cirkel snijdt de lijn, en de vlakken op de boog erboven glijden.
(b) Met $2\theta$ de hoek vanaf het $\sigma_1$-einde vraagt glijden
$45\sin2\theta \ge 0.85\,(65 + 45\cos2\theta)$, dat is
$45\sin2\theta - 38.25\cos2\theta \ge 55.25$, of
$59.06\sin(2\theta - \ang{40.36}) \ge 55.25$, waarbij $\ang{40.36} = \phi$.
Dus $\sin(2\theta - \phi) \ge 0.9355$, $2\theta - \phi$ tussen \ang{69.31}
en \ang{110.69}, $2\theta$ tussen \ang{109.67} en \ang{151.06}: de normalen
liggen op \ang{54.8} tot \ang{75.5} van $\sigma_1$, en de glijdende breuken
op \ang{14.5} tot \ang{35.2} van $\sigma_1$. Het midden van het bereik,
\ang{24.8}, is $\ang{45} - \phi/2$, het vlak van
\cref{thm:b1:stress-strain-rheology:failure-angle}. (c) Alleen dat vlak glijdt
wanneer de cirkel de lijn net raakt: $r = \mu(20 + r)/k$, $r = 20\mu/(k -
\mu) = 17/0.4624 = \qty{36.8}{MPa}$, en $\sigma_1' = 20 + 2r =
\qty{93.5}{MPa}$.
\end{solution}

\begin{solution}{pb:b1:stress-strain-rheology:1}
\textbf{1.} $\sigma_v = 2700\times9.81\times5000 = \qty{132.44}{MPa}$.

\textbf{2.} $P = 1000\times9.81\times5000 = \qty{49.05}{MPa}$, en
$\lambda = 49.05/132.44 = 0.370$.

\textbf{3.} $\sigma_1 = \sigma_v = \qty{132.44}{MPa}$, verticaal, en
$\sigma_3 = \qty{75.00}{MPa}$, horizontaal oost--west: de \omterm{def:b1:layered-earth:layers}{korst} wordt onder haar
eigen gewicht oost--west uitgerekt, de toestand die normaalbreuken aandrijft,
zoals deze.

\textbf{4.} $\sigma_1 - \sigma_3 = \qty{57.44}{MPa}$; middelpunt
$(132.44 + 75.00)/2 = \qty{103.72}{MPa}$; straal \qty{28.72}{MPa}.

\textbf{5.} $\sigma_1' = 132.44 - 49.05 = \qty{83.39}{MPa}$, $\sigma_3' =
75.00 - 49.05 = \qty{25.95}{MPa}$, middelpunt $103.72 - 49.05 = \qty{54.67}{MPa}$.

\textbf{6.} $\sigma_{xx} = \qty{75.00}{MPa}$, $\sigma_{yy} = \qty{132.44}{MPa}$
en $\sigma_{xy} = 0$: de assen zijn de hoofdrichtingen, waarvan er één verticaal
is, zodat de vlakken loodrecht erop geen schuifspanning dragen.

\textbf{7.} Langs de helling loopt de breuk langs de eenheidsvector
$(\cos\ang{60}, -\sin\ang{60})$, oostwaarts en omlaag; haar naar boven gerichte
normaal is $\vec n = (\sin\ang{60}, \cos\ang{60}) = (0.866, 0.500)$, onder
\ang{30} met $x$. Haar hoek met $\sigma_1$, langs $y$, heeft de cosinus
0.500: $\theta = \ang{60}$.

\textbf{8.} Met \ang{30} vanaf $x$: $\sigma_n = 75.00\times0.75 + 132.44\times0.25
= 56.25 + 33.11 = \qty{89.36}{MPa}$, en $\tau = (132.44 - 75.00)\times
\sin\ang{30}\cos\ang{30} = 57.44\times0.4330 = \qty{24.87}{MPa}$.

\textbf{9.} Bij $2\theta = \ang{120}$ vanaf het $\sigma_1$-einde:
$\sigma_n = 103.72 + 28.72\cos\ang{120} = 103.72 - 14.36 = \qty{89.36}{MPa}$
en $\tau = 28.72\sin\ang{120} = \qty{24.87}{MPa}$.

\textbf{10.} $\sigma_n' = 89.36 - 49.05 = \qty{40.31}{MPa}$.

\textbf{11.} $24.87/40.31 = 0.617$, onder 0.85: de breuk houdt stand.

\textbf{12.} $0.85\times40.31 = \qty{34.26}{MPa}$; de schuifspanning komt tekort
met \qty{9.39}{MPa}.

\textbf{13.} De cirkel heeft zijn middelpunt op \qty{54.67}{MPa} en de straal
\qty{28.72}{MPa}, van \num{25.95} tot \qty{83.39}{MPa}; de lijn van Byerlee is
$\tau = 0.85\,\sigma_n'$ (alle \omterm{def:b1:stress-strain-rheology:stress}{normaalspanningen} hier liggen onder
\qty{200}{MPa}); de breuk is het punt $(40.31, 24.87)$, bij \ang{120} vanaf het
$\sigma_1$-einde.

\textbf{14.} $0.85\times54.67/\sqrt{1 + 0.85^2} = 46.47/1.3124 =
\qty{35.41}{MPa}$, meer dan de straal, \qty{28.72}{MPa}: de hele cirkel ligt onder
de lijn, en geen vlak door de rots kan glijden.

\textbf{15.} $\phi = \arctan0.85 = \ang{40.36}$: het vlak op
$\ang{45} - \phi/2 = \ang{24.8}$ van de verticale $\sigma_1$, met een helling van
\ang{65.2}.

\textbf{16.} $S = \sigma_v' = \qty{83.39}{MPa}$ en
$\sigma_1 - \sigma_3 = 2\times0.85\times83.39/(1.3124 + 0.85) =
\qty{65.55}{MPa}$; de helling geeft $13.11\times5 = \qty{65.55}{MPa}$.

\textbf{17.} $\sigma_3 = 132.44 - 65.55 = \qty{66.88}{MPa}$.

\textbf{18.} De breuk glijdt wanneer
$r\sin\ang{120} = 0.85\,(c' + r\cos\ang{120})$ met
$c' = (83.39 + \sigma_3')/2$ en $r = (83.39 - \sigma_3')/2$; oplossen geeft
$\sigma_3' = \qty{17.18}{MPa}$, $\sigma_3 = \qty{66.23}{MPa}$. Dat is lager dan
voor het beste vlak: de breuk ligt niet op het punt waar de cirkel de lijn eerst
raakt, dus moet de cirkel verder groeien voordat hij het eigen punt van de breuk
bereikt.

\textbf{19.} Elke effectieve \omterm{def:b1:stress-strain-rheology:stress}{normaalspanning}, $\sigma_1'$ en $\sigma_3'$
inbegrepen, daalt met $\Delta P$, en hun verschil verandert niet; de
schuifspanningen hangen helemaal niet van de poriëndruk af.

\textbf{20.} Het beste vlak glijdt wanneer de afstand gelijk is aan de straal:
$0.85\,(54.67 - \Delta P)/1.3124 = 28.72$, dus
$\Delta P = 54.67 - 28.72\times1.3124/0.85 = 54.67 - 44.34 = \qty{10.33}{MPa}$.

\textbf{21.} De waterkolom in de put weegt evenveel als voorheen, dus moet de
pomp de volledige \qty{10.33}{MPa} aan de bronkop toevoegen, ongeveer 103 keer de
atmosferische druk: het gewicht van een waterkolom van
$\num{10.33e6}/(1000\times9.81) = \qty{1053}{m}$ hoog.

\textbf{22.} Met $\mu = 0.6$ kon de breuk slechts
$0.6\times40.31 = \qty{24.19}{MPa}$ dragen, minder dan de \qty{24.87}{MPa} die zij
draagt: zij zou al glijden. Omdat zij rustig is, is haar wrijving ten minste
$0.617$, of zij heeft een zekere \omterm{def:b1:stress-strain-rheology:coulomb}{cohesie}.

\textbf{23.} Het best georiënteerde vlak glijdt wanneer de poriëndruk bereikt
$49.05 + 10.33 = \qty{59.38}{MPa}$, dat is $\lambda = 59.38/132.44 = 0.448$,
ruim onder de lithostatische waarde 1.

\textbf{24.} $\Delta P = \sigma_n' - \tau/0.85 = 40.309 - 24.870/0.85 = 40.309 -
29.259 = \qty{11.049}{MPa}$, meer dan de \qty{10.33}{MPa} die een vlak van de beste
oriëntatie tot schuiven zou brengen. Op één decimaal is de stijging van de
poriëndruk die de breuk tot schuiven brengt
\textbf{\qty[text-series-to-math=true]{11.0}{MPa}}.
\end{solution}
